\subsection{Positive Hermitian forms.}

DEFINITION 1. --- A Hermitian form \(\Phi\) on E is said to be positive (resp. negative) if \(\Phi(x, x) \geqslant 0\) (resp. \(\Phi(x, x) \leqslant 0\) ) for all \(x \in \mathbf{E}\) .

When A = K, one also says that the quadratic form \(Q(x) = \frac{1}{2} \Phi(x, x)\) to which \(\Phi\) is associated is positive (resp. negative).

Suppose E is finite-dimensional over A, and let \((e_{i})\) ( \(i = 1, \ldots, n\) ) be an orthogonal basis of E ( \(\S 6, no. 1\) , th. 1). For a Hermitian form \(\Phi\) on E to be positive, it is necessary and sufficient that \(\Phi(e_{i}, e_{i}) \geqslant 0\) for \(i = 1, \ldots, n\) . Let \(\Phi\) be a nondegenerate positive Hermitian form on E; since every positive element of K is a square, hence of the form \(\rho\bar{\rho} (\rho \in A)\) , there exist in E orthonormal bases for \(\Phi\) ( \(\S 6, no. 1\) , cor. 1 of th. 1) .

PROPOSITION 1. --- Suppose E is finite-dimensional, and A = K or A = K(i). If \(\Phi\) is a nondegenerate positive Hermitian form on E, then the extensions of \(\Phi\) to \(\bigotimes^{p}\mathbf{E}\) and \(\bigwedge^{p}\mathbf{E}\) ( \(p > 0\) ), as well as the inverse form of \(\Phi\) , are nondegenerate positive Hermitian forms.

This follows immediately from the existence of an orthonormal basis of E, and from prop. 2, § 6, no. 1.

Prop. 1 remains true if one replaces everywhere “nondegenerate positive” by “positive”.

PROPOSITION 2. --- Let Φ be a positive Hermitian form on E. For x, y in E, one has

\[
\Phi (x, y) \overline {{{{\Phi (x , y)}}}} \leqslant \Phi (x, x) \Phi (y, y).\tag{1}
\]

The inequality is indeed immediate when the vectors \(x\) and \(y\) are proportional. Suppose therefore \(x\) and \(y\) linearly independent. Let A' be the (commutative) subfield K( \(\Phi(x, y)\) ) of A, F the vector plane over A' generated by \(x\) and \(y\) , and \(\Phi_{\mathbb{F}}\) the restriction of \(\Phi\) to F; this takes its values in A'. By prop. 1, the discriminant of \(\Phi_{\mathbb{F}}\) with respect to the basis \((x, y)\) is \(\geqslant 0\) . Now this discriminant is \(\Phi(x, x)\Phi(y, y) - \Phi(x, y)\overline{\Phi(x, y)}\) . QED.

COROLLARY. --- The set of isotropic vectors of E is the subspace \(E^{0}\) orthogonal to E for \(\Phi\) . For \(\Phi\) to be nondegenerate, it is necessary and sufficient that one have \(\Phi(x, x) > 0\) for all \(x \neq 0\) .

PROPOSITION 3. --- Suppose E is finite-dimensional and A is commutative (A = K or A = K(i)). Let \(\Phi\) be a Hermitian form on E, X its matrix with respect to a basis ( \(x_j\) ) ( \(j = 1, \ldots, n\) ) of E; for every subset H of {[}1, n{]}, denote by \(X_{\mathrm{H,H}}\) the minor of X obtained by deleting the rows and columns with indices \(j \notin \mathrm{H}\) (chap.~III, § 6, no. 3).

\begin{enumerate}
\def\labelenumi{\alph{enumi})}
\item
  If \(\Phi\) is nondegenerate positive, one has \(X_{\mathrm{H,H}} > 0\) for every subset H of \([1, n]\) .
\item
  Conversely, if, setting \(H_{j} = (1, j)\) , one has \(X_{H_{j}, H_{j}} > 0\) for \(j = 1, \ldots, n\) , \(\Phi\) is nondegenerate positive.
\end{enumerate}

Suppose first \(\Phi\) nondegenerate positive. The elements \((x_{j})\) ( \(j\in\mathbb{H}\) ) form a basis of a subspace F of E, and the minor \(X_{\mathrm{H,H}}\) is the discriminant of the restriction \(\Phi_{\mathrm{F}}\) of \(\Phi\) to F with respect to this basis; now, since \(\Phi_{\mathrm{F}}(x,x)>0\) for all \(x\neq0\) in F, \(\Phi_{\mathrm{F}}\) is nondegenerate positive (cor. of prop. 2); one therefore has \(X_{\mathrm{H,H}}>0\) (prop. 1). To prove \(b\) ), note that, with the notations of prop. 1, § 6, no. 1, the minor \(X_{\mathrm{Hj,Hj}}\) is equal to D \(_{j+1,j+1}\) ; there therefore exists (prop. 1, § 6, no. 1) an orthogonal basis \((e_{j})\) ( \(j=1,\ldots,n\) ) of E such that \(\Phi(e_{j},e_{j})>0\) for \(j=1,\ldots,n\) ; consequently \(\Phi\) is nondegenerate positive.

Remark. --- It follows from the existence of orthonormal bases that two nondegenerate positive Hermitian forms on two vector spaces of the same finite dimension are equivalent (§ 1, no. 6). Let then L be a finite-dimensional Hermitian space over A, whose metric form is nondegenerate positive, and V₁ and V₂ two linear varieties of the same dimension in L (§ 6, no. 6); since the restrictions of the metric form to the directions T₁ and T₂ of V₁ and V₂ on the one hand, and to the orthogonal subspaces T₁ and T₂ on the other hand, are equivalent, there exists a unitary automorphism \(u\) of the space T of translations of L such that \(u(\mathrm{T}_1) = \mathrm{T}_2\) and \(u(\mathrm{T}_1^0) = \mathrm{T}_2^0\) ; there therefore exists a displacement \(v\) of L such that \(v(\mathrm{V}_1) = \mathrm{V}_2\) . Let \((a, b), (a', b')\) two pairs of points of L; for there to exist a displacement \(v\) of L such that \(v(a) = a'\) and \(v(b) = b'\) it is necessary and sufficient (with the notation of § 6, no. 6) that one have \(e(a, b) = e(a', b')\) ; the element \(\sqrt{e(a, b)}\) of A (chap.~VI, § 2, no. 4) is called the distance of \(a\) and \(b\) in the Hermitian space L.

